(1) 直接有 \(\Phi_{xx}=(U+V\cos\Omega t)/r_0^2\) 与 \(\Phi_{yy}=-(U+V\cos\Omega t)/r_0^2\)，两者抵消，满足 \(\nabla^2\Phi=0\)。Newton方程为
\[
\ddot x+\frac{e(U+V\cos\Omega t)}{mr_0^2}x=0,\qquad
\ddot y-\frac{e(U+V\cos\Omega t)}{mr_0^2}y=0.
\]
写 \(\tau=\Omega t/2\)，每一次时间求导产生 \(\Omega/2\)，故标准Mathieu参数是
\[
\boxed{a_x=\frac{4eU}{mr_0^2\Omega^2},\quad q_x=-\frac{2eV}{mr_0^2\Omega^2},\qquad
a_y=-a_x,\quad q_y=-q_x}.
\]
当 \(V=0\) 且 \(U\ne0\)，两个方向的回复系数符号相反，至少一轴指数失稳。没有静态平面四极电势可以两轴同时稳定；动态周期驱动是克服这一局部鞍点障碍的关键，不能把电场线曲率错误当作等向回复力。

(2) 记 \(\epsilon=eV/(mr_0^2)\)。取 \(x=X+\xi_x,\ y=Y+\xi_y\)，\(X,Y\) 比射频周期变化慢。消去一阶振荡力的快解为
\[
\xi_x=\frac{\epsilon}{\Omega^2}X\cos\Omega t,\qquad
\xi_y=-\frac{\epsilon}{\Omega^2}Y\cos\Omega t.
\]
回代方程并使用周期平均 \(\langle\cos^2\Omega t\rangle=\tfrac12\)，得
\[
\ddot X+\frac{\epsilon^2}{2\Omega^2}X=0,\qquad
\ddot Y+\frac{\epsilon^2}{2\Omega^2}Y=0.
\]
于是 \(\boxed{\omega_{\rm sec}=|\epsilon|/(\sqrt2\,\Omega)}\)，
\[
\boxed{U_{\rm ps}=\frac{m\epsilon^2}{4\Omega^2}(X^2+Y^2)
=\frac{e^2V^2}{4m\Omega^2r_0^4}(X^2+Y^2)}.
\]
由于 \([\epsilon]=T^{-2}\)，上述频率及势能量纲正确；符号 \(e\to-e\) 仅反转快微运动相位，赝势仍为正。需要 \(|\epsilon|/\Omega^2\ll1\) 且世俗频率远小于射频频率，不能在失稳舌边缘滥用平均。

(3) 当 \(X,Y\) 近乎冻结时，快速度为
\(\dot\xi_x=-\epsilon X\sin\Omega t/\Omega\)，
\(\dot\xi_y=\epsilon Y\sin\Omega t/\Omega\)。故
\[
\langle K_{\rm micro}\rangle=
\frac{m\epsilon^2}{4\Omega^2}(X^2+Y^2).
\]
射频电势能为 \(e\Phi_{\rm rf}=m\epsilon(x^2-y^2)\cos\Omega t/2\)。把 \(x=X+\xi_x,\ y=Y+\xi_y\) 展开，并消去奇数周期平均，得
\[
\langle e\Phi_{\rm rf}\rangle=
\frac{m\epsilon^2}{2\Omega^2}(X^2+Y^2)
=2\langle K_{\rm micro}\rangle.
\]
因此是 \(\langle K_{\rm micro}-e\Phi_{\rm rf}\rangle=-U_{\rm ps}\)，而不是简单把两种正能量相加来定义赝势。真实系统
\(H=\tfrac12m(\dot x^2+\dot y^2)+e\Phi\) 显含时间，沿Newton轨道有
\[
\boxed{\dot H=e\partial_t\Phi=-\frac{eV\Omega}{2r_0^2}(x^2-y^2)\sin\Omega t}.
\]
电源可以瞬时输入或抽取能量；慢变量有效能量的近似守恒不等于原粒子机械能瞬时守恒。

(4) 直流势在平均方程中给 \(\pm eU/(mr_0^2)\) 的曲率，故
\[
\boxed{\omega_x^2=\frac{\epsilon^2}{2\Omega^2}+\frac{eU}{mr_0^2}},\quad
\boxed{\omega_y^2=\frac{\epsilon^2}{2\Omega^2}-\frac{eU}{mr_0^2}}.
\]
两轴平均线性束缚要求
\[
\boxed{\left|\frac{eU}{mr_0^2}\right|<
\frac{e^2V^2}{2m^2\Omega^2r_0^4}}.
\]
等号时一轴的慢振动软化到零；零频不等于有恢复力的稳定束缚。\(V\to0\) 时赝势消失，非零 \(U\) 形成静电鞍点；若 \(U=V=0\) 则是中性自由漂移而非束缚。固定 \(U,V\) 而 \(\Omega\to\infty\)，射频平均约束趋零，微运动消失，直流鞍点依然存在；实际稳定图在较大Mathieu参数处有额外的Floquet共振与稳定舌，需使用单周期传递矩阵精确判断。